\documentclass[12pt, a4paper]{ctexart}
\usepackage{geometry}
\geometry{left=2.5cm,right=2.5cm,top=2.5cm,bottom=2.5cm}
\usepackage{hyperref}
\usepackage{enumitem}

\begin{document}

\begin{titlepage}
    \centering
    \vspace*{2cm}
    
    {\Huge \bfseries 软件工程与实践课程项目 \par}
    \vspace{1.5cm}
    
    {\Huge \bfseries 可行性分析报告 \par}
    \vspace{2cm}
    
    {\Huge \bfseries 项目名称：Mini Blog \par}
    \vspace{4cm}
    
    \Large
    \renewcommand{\arraystretch}{1.5}
    \begin{tabular}{|c|c|c|}
        \hline
        \textbf{小组成员} & \textbf{姓名} & \textbf{学号} \\
        \hline
        开发人员 & 刘俊汝 & 202300172027 \\
        \hline
    \end{tabular}
    
    \vfill
    
    {\Large 2026年4月 \par}
\end{titlepage}

\setcounter{page}{1} % 让目录之后的内容从第1页开始
\tableofcontents
\newpage

\section{引言}
\subsection{目的与标识}
本文档是博客网站系统“Mini Blog”项目的可行性分析报告，主要从业务需求、技术实现、经济成本及法律合规等方面对该项目进行综合评估，为下一步的需求分析和系统设计提供决策依据。

\subsection{项目背景}
本项目是在软件工程实践课程要求下提出的。当前互联网内容分享高度发达，建立一个集内容产出、展示与互动交流于一体的平台具有充分的实际意义。在此项目背景下，旨在通过实践掌握软件项目的完整生命周期管理体系与开发技术。

\subsection{项目概述}
本系统旨在构建一个轻量级、交互友好的个人与公共博客平台。系统涉及两种主要角色：注册用户和系统管理员。涵盖的核心功能模块包括：
\begin{itemize}
    \item \textbf{用户注册及登录模块}：实现用户的身份认证，验证通过后进入个人博客页面。
    \item \textbf{用户及功能模块}：用户能够自由发表、修改和删除个人的文章，并在浏览阶段支持对个人及其它作者博文的评论交流。
    \item \textbf{后台管理模块}：面向系统管理员，用于对网站的基础信息、用户信息、文章及评论内容进行集中化的修改、查询与删除处理。
\end{itemize}

\section{可行性分析的前提}
\subsection{项目的要求与目标}
本项目旨在课程规定的时限内，开发出一套功能完备、运行稳定的B/S（Browser/Server）架构Web应用系统。系统需具备清晰直观的UI界面和健壮无误的核心业务逻辑。

\subsection{环境、条件与限制}
作为一门课程的实践项目，开发周期较短，系统主要在校园环境或个人计算机中完成搭建、编码与测试。因为没有专门的项目研制资金支持，所以全部的架构设计与实现均需依托开源免费的技术栈、框架以及数据库系统。

\section{可选的方案}
针对博客系统这类偏重图文展示与多端访问的应用，目前可选的主要客户端方案包括传统的 Desktop C/S（客户端/服务器）系统和基于 Web 的 B/S 系统。

\subsection{可选择的系统方案1：C/S 架构方案}
基于C/S（Client/Server，客户端/服务端）架构的研究方案。该方案要求用户在PC端安装专用的桌面客户端软件（例如使用C\# WinForms、WPF或Java Swing等技术开发）。桌面客户端通过内部网络或互联网，直接与数据库或应用服务器进行数据通信。此方案的优势在于可以充分利用本地系统资源，软件运行效率高，能够设计丰富的原生图形界面；但劣势非常明显，即用户需要手动下载、安装且频繁更新客户端。此外，不同的操作系统（Windows、macOS、Linux）难以做到一套代码跨平台完美兼容。对用户而言访问成本高、传播分享困难，完全不匹配博客网站天然互联网轻量级分享与阅读的特性。

\subsection{可选择的系统方案2：B/S 架构方案}
基于B/S（Browser/Server，浏览器/服务器）架构的Web应用方案。用户无需安装任何专门的客户端软件，只要通过操作系统内建的Web浏览器（如Chrome、Edge、Safari等）直接访问目标网站URL即可进入系统。其开发模式可采用主流的前后端分离体系（例如后端采用 Python Django 提供 RESTful API 或模板渲染服务，前端使用原生的 HTML、CSS、JavaScript 来构建动态页面）。这种模式具有天然的强跨平台特性，应用部署与更新均由服务端单方面完成，用户端随时享有最新版本功能；同时内容的分享只需发送一个网页链接，极大地便利了内容的传播，与现代读者消费博客及参与评论社交的场景完美吻合。

\subsection{选择最终方案的准则}
根据课程实践要求与所立课题的特性，最终方案的选取遵循以下几项核心准则：
\begin{enumerate}
    \item \textbf{跨平台能力与易用性}：系统必须支持各种主流操作系统的随时访问，保障用户使用门槛尽可能低，能够随时随地进行阅读、发表与互动。
    \item \textbf{开发成本与实现周期}：由于项目受限于课程较短的开发时间，决定的方案应具备高度成熟的技术生态与开源可复用的框架，从而降低系统开发难度，提高研发测试效率。
    \item \textbf{维护更新透明化}：后续可能的系统系统升级与修复（如管理员后台增补功能）不需要用户在终端设备上进行任何干预，直接由服务端进行热更新，实现用户无感知升级。
    \item \textbf{内容的社交分享属性}：文章页面应具备直接的入口链接分享机制，极大方便用户通过各类即时通讯软件发送并引流。
\end{enumerate}
基于以上准则评估，方案1在跨平台推广属性、维护便捷度以及开发效率上全面落后于方案2。因此，综合评价后，本系统采用 \textbf{系统方案2（B/S 架构方案）} 作为最终架构选型。

\section{所建议的系统}
\subsection{对所建议的系统的说明}
基于前述架构选型论证，所建议的博客系统（Mini Blog）采用 B/S 架构搭建，以保障代码的高内聚与低耦合。系统采用了经典的三层设计模式：
\begin{itemize}
    \item \textbf{表示层（Presentation Layer）}：采用原生的 HTML5、CSS3 及 JavaScript 技术构建前端响应式展示视图与交互表单。
    \item \textbf{业务逻辑层（Business Logic Layer）}：依托 Python Django 框架为核心，处理 HTTP 请求指令，统一调度各种功能业务逻辑的流转。
    \item \textbf{数据持久层（Data Access Layer）}：利用 Django 强健的内置 ORM 机制，将底层 MySQL 8.0 数据库的 SQL 操作转化为安全的面向对象交互。
\end{itemize}

\subsection{数据流程和处理流程}
对于Mini Blog这套博客网站系统的核心数据流向与处理，可抽象归纳为以下几条主线：
\begin{enumerate}
    \item \textbf{用户注册与认证流程}：
    访客在前端页面填写注册信息（账户名、加密凭证等），前端将请求载体发送到服务端；后端校验字段并在数据库用户表中验证。若信息合法即生成记录并持久化密码；登录时后端比对信息，匹配成功则生成并下发会话控制（Session）以用于维持后续身份层面的合法性。
    \item \textbf{文章管理流程（增删改查数据流）}：
    作者通过前端编辑器提交包含标题及正文的文章数据。数据经逻辑规则验证及重组装后，通过服务器 ORM 层直接落库保存到数据表中。随即页面触发更新查询指令，向服务端索取最新的文章目录集合进而渲染更新。
    \item \textbf{评论互动处理流程}：
    阅读者经由页面将填写的文本连带被评论文章的唯一标记发送至后端 API；后端封装该部分数据存入评论数据表，通过外键约束挂载并归属于对应的文章下。
    \item \textbf{平台数据审批及处理}：
    系统管理员以高级权限账号认证后，可通过独立接口获取全站业务数据集合（用户信息、博文及其相关评论等）。经拦截鉴别器放行后，管理员可对数据表发送越权覆盖动作，实现统一的信息管控流。
\end{enumerate}

\subsection{影响与要求}
\subsubsection{设备}
普通的个人电脑即可支撑其研发及前中期的代码运行、局域网内预览。在验收和上线展示部署期，只需提供一台入门级的云服务器用以实现公网承载。

\subsubsection{软件}
使用目前开源主流的 Web 技术堆栈。运行环境的操作系统不受限（开发端使用 Windows11），开发集成环境使用 PyCharm 和 VS Code。服务端业务核心依托 Python 语言，数据库服务器要求安装有 MySQL 8.0 版服务。

\subsubsection{运行}
因基于 B/S 架构，系统的运行依托云端的主机守护。只需在服务端采用 Nginx 或者 WSGI 应用容器长期运行驻留 Web 服务，且维持 MySQL 连通即可。用户一端真正做到了免配置，仅需标准 Web 浏览器即可运行操作。

\subsubsection{开发}
作为短周期的课程实践课题项目，由组员在特定的时间节点内单兵作战即可。

\subsubsection{环境}
编码开发环境只需利用校园局域网络搭配开源工具包；而若是正式部署，环境上需注意设定主机的网络安全策略和安全组防火墙策略，如开启 80、443 协议入口。

\subsubsection{经费}
系统的设计、框架选型以及核心语言解释器环境都是基于开源免费组件搭建，不涉及商业版权的纠纷或购买负担；唯一可能的微乎其微的费用是上线所租赁的短时测试型云主机服务费，综合研发成本基本为零。

\subsection{局限性}
此方案作为教学性质项目立项完成，主要在如下存在局限性：
首先，业务模型的设计难以满足企业级生产线上大规模的高并发吞吐需求，只可用于轻量级别的演示和小圈子内容生产；其次，界面交互可能较为缺乏细致深入的动效打造，在感官上达不到大厂级的工业标准。

\section{经济可行性分析}
本项目为教学实践项目，不涉及商业投资利润目标。
\begin{itemize}
    \item \textbf{投资总额}：开发所使用的IDE软件、框架、数据库系统（如MySQL）等均采用免费开源软件版本，且不需要商用服务器与专职的人工成本薪酬，整体开发资金投入约等于 0 元。
    \item \textbf{预期收益}：经济收益虽无法量化，但项目能够为组员提供极高的隐藏“非一次性收益”——切实提高成员的软件工程规范认知水平和编程实践能力，顺利通过课程考核。
\end{itemize}
\textbf{结论}：经济上具有绝佳的可行性。

\section{技术可行性分析}
博客系统具备非常典型的系统特征，以信息库的“增删改查”（CRUD操作）为核心。
\begin{itemize}
    \item 互联网开源社区上有着极其丰富且成熟的源码、架构解析与开发文档可供借鉴。
    \item 选用的主开发语言和数据库技术方案，项目成员在前期专业课程中已做过一定深度的入门与掌握。
    \item 当下各个Web框架的功能封装集成度极高，开发难度大为降低。
\end{itemize}
\textbf{结论}：该项目在当前人员与技术储备条件下具备充足的技术可行性。

\section{法律可行性}
\begin{itemize}
    \item \textbf{侵权与违法责任}：本系统从零开始编码设计，或在遵循且保护原开源许可协议（如MIT/Apache）前提下引用第三方库。开发全流程属于学术实践学习范畴，不构成抄袭与商业知识产权侵犯。系统不包含涉及国家安全保密或法律严令禁止的内容功能。
\end{itemize}
\textbf{结论}：法律上完全可行无风险。

\section{用户使用可行性}
博客网站的大众普及率极高，普通用户和管理员只要具备基础的互联网浏览器操作常识，就能非常直觉化地完成登录、操作发布与评论等功能交互。无需对用户提供复杂的使用培训。
\textbf{结论}：用户操作可行性极高。

\section{综合结论}
综合对当前项目技术、经济、法律风险及用户使用的系统评估，该课题定位恰当，功能拆解明确，各方面资源符合软件工程实践课程背景。本项目完全具备实施价值与可行依据，具备立项条件。

\end{document}